\section{Prayer and the Enlargement of Love}

We can speak rightly about our final end and still live with our attention almost entirely elsewhere. Prayer turns what we know of God into personal address and consent. We attend to the one who knows and loves us, bringing him the life we hope he will perfect. The act may be simple: listening to the Gospel, recollecting ourselves before a task, remaining quietly with Christ, or giving thanks for a received good. Through such acts we learn to seek God during a day otherwise occupied with his gifts.

Bernard begins \emph{On the Love of God} by saying that God himself is the reason for loving God (I). He unfolds the claim through the prior generosity of divine love. God has no need of what we can supply; he has loved sinners and given his Son for them. Our love answers his prior gift. This prevents the desire for heaven from becoming a negotiation in which God is valued only as a means to some other reward. He gives himself. To desire that gift rightly is to desire him, and the growth of charity makes our desire more truthful.

This growth does not require us to stop wanting happiness. God made us for blessedness, and the hope of receiving it is good. What needs purification is the separation we sometimes make between our happiness and God's goodness. We can treat prayer as a means of obtaining control, comfort, or exemption from the claims of love. Bernard directs us to the God whose own love is free and generous. As we come to know him, we can desire our good within his will and learn to rejoice in what honors him. The promised reward and the divine giver are united because the reward is communion with him.

In his account of the highest degree of love, Bernard speaks of loving oneself for God's sake (X). He describes consent to God's will extending even to the love we bear ourselves. His images of water mingled with wine, iron glowing in fire, and air filled with light suggest a creature thoroughly changed by participation in divine love. He then states that human nature is preserved and receives a greater beauty, power, and glory. We remain the persons God created while learning to love ourselves in relation to him.

To love ourselves in God is also to receive our limitations without hatred. Some forms of apparent humility conceal resentment that we are not more impressive, less needy, or less dependent upon mercy. True humility can ask for growth while giving thanks for existence itself. God does not wait for us to become self-sufficient before he can love us. Prayer gradually brings this truth into the heart. We can acknowledge a fault without turning against the whole person who needs forgiveness. We can accept help, submit a talent to service, and allow God to choose forms of fruitfulness that do not flatter our ambitions.

The same purification reaches prayer's consolations. A period of sweetness can awaken gratitude and strengthen perseverance, yet the heart may begin to seek that experience as though it were the end. When the feeling passes, we can discover whether we are willing to remain with the Lord. Faithful prayer in such a season does not prove that love has become perfect, but it gives love room to become less possessive. We return to the received words of Scripture and the Church, offer our attention as we can, and ask for the charity that makes prayer true even when it is poor.

Francis de Sales's meditation on paradise allows the imagination to serve this desire. He takes the beauty of a starlit night and a brilliant summer day, then points beyond their combined splendor to the blessed country (I.16). He moves from beauty to the saints and from the saints to God's presence, where their joy is fulfilled. Our own meditation can begin with something beautiful we have seen and received gratefully. In giving thanks, we turn to its Creator and can come to desire the uncreated good from whom its loveliness comes.

The person who desires God's presence must also receive God's will in the place where he has been set. A request that interrupts our recollection may call for the patience we have just sought in prayer. Our reluctance shows how much we still need the grace for which we have asked. Answering the request gives us an occasion to act upon that grace and to love the neighbor whose need has become apparent.

